\section{The sublinear range and a formula at every order}\label{sec:asymptotics}

The square-root correction is the first term of a convergent exponent. Retaining that exponent describes every sublinear order. Keeping the permutation denominator exact then gives one formula throughout the full range $2\le k\le N$.

Define, for $0\le t<1$,
\begin{equation}\label{eq:actionG}
G(t)=t+\frac{t^2}{2}+(1-t^2)\log(1-t)-\frac23t^3
=\sum_{j\ge4}\frac{2t^j}{j(j-2)}.
\end{equation}
All coefficients in the last series are positive.

\begin{theorem}[Uniform sublinear approximation]\label{thm:sublinear}
For every sequence $b_N=o(N)$ with $b_N\ge2$,
\begin{equation}\label{eq:sublinear}
\sup_{2\le k\le b_N}
\left|\frac{R_k}{H_k\exp\{NG(\sqrt{k/N})\}}-1\right|
\longrightarrow0.
\end{equation}
The supremum is over integers. The definition of $H_k$ remains \eqref{eq:hyperdef}, including $K=k-1$ at odd orders.
\end{theorem}

\begin{corollary}[Sharp thresholds for successive corrections]\label{cor:hierarchy}
For a fixed integer $p\ge4$, put
\[
G_p(t)=\sum_{j=4}^{p}\frac{2t^j}{j(j-2)}.
\]
Along any sequence $2\le k=o(N)$,
\begin{equation}\label{eq:hierarchy}
\frac{R_k}{H_ke^{NG_p(\sqrt{k/N})}}\longrightarrow1
\quad\Longleftrightarrow\quad
k=o\left(N^{(p-1)/(p+1)}\right).
\end{equation}
The ratio stays bounded exactly when $k^{p+1}/N^{p-1}$ stays bounded. At $k/N^{(p-1)/(p+1)}\to c\in(0,\infty)$ its limit is
\begin{equation}\label{eq:hierarchyboundary}
\exp\left(\frac{2c^{(p+1)/2}}{(p+1)(p-1)}\right).
\end{equation}
In particular, the correction $e^{k^2/(4N)}$ is relatively accurate up to $k=o(N^{3/5})$. Adding $2k^{5/2}/(15N^{3/2})$ to its exponent extends that range to $k=o(N^{2/3})$.
\end{corollary}

The thresholds in \eqref{eq:hierarchy} are uniform upper-range statements as well: the relative error tends uniformly to zero on $2\le k\le b_N$ whenever $b_N=o(N^{(p-1)/(p+1)})$. The conclusion concerns each fixed truncation degree $p$.

\subsection{A coefficient estimate with both angular peaks controlled}

Write
\[
f(z)=\frac{z}{1-z}+\log(1-z),\qquad
x=\sqrt{K/N},\qquad r=\frac{x}{1+x}.
\]
The positive saddle and its variance parameter are explicit:
\begin{equation}\label{eq:saddleequation}
Nrf'(r)=\frac{Nr^2}{(1-r)^2}=K,\qquad
b=N(r\partial_r)^2f(r)=2K(1+x).
\end{equation}

\begin{lemma}[Uniform positive-saddle estimate]\label{lem:saddle}
Let $2\le K\le N$ be integers, $\epsilon\in\{0,1\}$ and $K+\epsilon\le N$. Put $\lambda=(2/3)Nx^3$. There is an absolute constant $C$ such that
\begin{equation}\label{eq:saddlecoefficient}
[z^{K+\epsilon}]e^{Nf(z)}
=\frac{e^{Nf(r)}r^{-K-\epsilon}}{\sqrt{2\pi b}}(1+E),
\qquad |E|\le C\left(K^{-1}+e^{-\lambda/8}\right).
\end{equation}
The right side is a relative approximation when $K$ and $\lambda$ diverge. The bound remains valid without that divergence.
\end{lemma}
\begin{proof}
Cauchy's formula reduces the normalized coefficient to
\[
\frac1{2\pi}\int_{-\pi}^{\pi}
e^{\Phi(\theta)-i\epsilon\theta}\,d\theta,
\qquad
\Phi(\theta)=N(f(re^{i\theta})-f(r))-iK\theta.
\]
The positive coefficients of $f$ give
\begin{equation}\label{eq:angulardamping}
-\Re\Phi(\theta)
\ge Nr^2\sin^2\theta+\frac23Nr^3(1-\cos3\theta).
\end{equation}
Here $r\le1/2$ and $Nr^2=K/(1+x)^2\ge K/4$. On $|\theta|\le\pi/6$,
\[
\Re\Phi(\theta)\le-\frac9{4\pi^2}K\theta^2.
\]
For each fixed $q\ge2$, the derivative sum
\[
N\sum_{j\ge2}\frac{j-1}{j}j^qr^j
\le K\sum_{j\ge2}\frac{j-1}{j}j^q2^{2-j}
\]
is bounded by an absolute constant depending only on $q$, times $K$. Taylor's inequalities and \eqref{eq:saddleequation} therefore give
\[
\Re\Phi=-b\theta^2/2+O(K\theta^4),\qquad
|\Im\Phi|\le CK|\theta|^3.
\]
Pair conjugate angles. The central integrand becomes
$e^{\Re\Phi}\cos(\Im\Phi-\epsilon\theta)$. Using
$|\cos u-1|\le u^2/2$ and the real exponential mean-value bound, its difference from $e^{-b\theta^2/2}$ is at most
\[
C(K\theta^4+K^2\theta^6+\theta^2)e^{-cK\theta^2}.
\]
Its integral is $O(K^{-3/2})$, whereas the Gaussian integral is of order $K^{-1/2}$. The Gaussian tails outside the central interval are exponentially small.

The pair term alone does not suppress the other peak. Write $\theta=\pi+u$ modulo $2\pi$, with $|u|\le\pi/6$. Then
\[
1-\cos3\theta=1+\cos3u\ge1,\qquad
\frac23Nr^3=\frac\lambda{(1+x)^3}\ge\frac\lambda8.
\]
The pair term in \eqref{eq:angulardamping} supplies Gaussian decay in $u$. The second-peak integral is therefore at most $CK^{-1/2}e^{-\lambda/8}$. At the remaining angles, $|\sin\theta|\ge1/2$ and the integral is at most $Ce^{-K/16}$. Dividing by the central Gaussian integral proves \eqref{eq:saddlecoefficient}. The circle $r\le1/2$ stays away from the singularity at one, including at $K=N$.
\end{proof}

\subsection{Restoring parity and the permutation denominator}

For the moment normalization take $K=2\lfloor k/2\rfloor=2m$ and $\epsilon=k-K$. Set
\begin{equation}\label{eq:actionq}
h(x)=x+(x^2-1)\log(1+x)-\frac{x^2}{2},\qquad
q(x)=h(x)-\frac23x^3.
\end{equation}
The action identity is exact:
\[
Nf(r)-K\log r+K\log x-K/2=Nh(x).
\]
Stirling's formula applied to $m!$ in $\mathcal L_k$, followed by Lemma~\ref{lem:saddle}, gives
\begin{equation}\label{eq:normalizedsaddle}
\frac{Q_k}{\mathcal L_k}
=\frac{e^{Nh(x)}(1+x)^{\epsilon-1/2}}{2\lambda^\epsilon}
\left[1+O\left(K^{-1}+e^{-\lambda/8}\right)\right].
\end{equation}
Indeed, the even factorial cancellation leaves $2^mm!N^{-m}$ times the coefficient. At odd orders the additional factor is exactly
\[
\frac3{2Kr}=\frac{1+x}{\lambda}.
\]
Thus no odd-order shift is lost. The exact hyperbolic term is
$H_k=(e^\lambda+(-1)^\epsilon e^{-\lambda})/(2\lambda^\epsilon)$.
When $K$ and $\lambda$ exceed sufficiently large fixed constants, taking logarithms in \eqref{eq:normalizedsaddle} yields
\begin{equation}\label{eq:logsaddle}
\log\frac{R_k}{H_k}
=\log D_k+Nq(x)+(\epsilon-1/2)\log(1+x)
+O\left(K^{-1}+e^{-\lambda/8}\right).
\end{equation}

For the sublinear formula, put $t=\sqrt{k/N}$ and
$I(u)=u+(1-u)\log(1-u)$. Monotonicity of $-\log(1-v)$ bounds its left Riemann sum by
\begin{equation}\label{eq:denominatorintegral}
0\le NI(k/N)-\log D_k\le-\log(1-k/N).
\end{equation}
The derivative identities
$h'(x)=2x\log(1+x)$ and $q'(x)=2x(\log(1+x)-x)$ give $|q'(x)|\le x^3$. Since $t^2-x^2=\epsilon/N$,
$N|q(t)-q(x)|\le k/N$. Finally, $I(t^2)+q(t)=G(t)$ exactly. Hence, uniformly for $k/N\le1/16$ under the same large-$K$, large-$\lambda$ conditions,
\begin{equation}\label{eq:logsublinear}
\log\frac{R_k}{H_k}
=NG(\sqrt{k/N})+(\epsilon-1/2)\log(1+\sqrt{K/N})
+O\left(k/N+K^{-1}+e^{-\lambda/8}\right).
\end{equation}
The parity term has order $\sqrt{k/N}$ and must be retained for finer error estimates.

\begin{proof}[Proof of Theorem~\ref{thm:sublinear}]
On $\sqrt N<k\le b_N=o(N)$, the minimum $K$ is of order $\sqrt N$ and the minimum $\lambda$ is of order $N^{1/4}$. Every error in \eqref{eq:logsublinear}, and its explicit parity term, tends uniformly to zero. On $2\le k\le\sqrt N$, Theorem~\ref{thm:sqrt} applies, while the positive series \eqref{eq:actionG} gives
\[
0\le NG(\sqrt{k/N})-\frac{k^2}{4N}
\le C\frac{k^{5/2}}{N^{3/2}}\le CN^{-1/4}.
\]
This patches the two ranges and includes fixed orders. It does not assume a single-saddle relative approximation when $\lambda$ is small.
\end{proof}

\begin{proof}[Proof of Corollary~\ref{cor:hierarchy}]
For fixed $p$, positivity and decreasing coefficients give
\[
G(t)-G_p(t)=\frac{2t^{p+1}}{(p+1)(p-1)}(1+O_p(t))
\quad(t\to0).
\]
The first term bounds the tail from below, and the same term divided by $1-t$ bounds it from above. Theorem~\ref{thm:sublinear} therefore gives
\[
\log\frac{R_k}{H_ke^{NG_p(\sqrt{k/N})}}
=\frac{2k^{(p+1)/2}}{(p+1)(p-1)N^{(p-1)/2}}
\bigl(1+O_p(\sqrt{k/N})\bigr)+o(1).
\]
This proves relative accuracy, its necessity, boundedness and the boundary limit. The same positive-tail bound proves uniformity over the stated upper ranges. In particular,
\[
NG(\sqrt{k/N})
=\frac{k^2}{4N}+\frac{2k^{5/2}}{15N^{3/2}}
+\frac{k^3}{12N^2}+\frac{2k^{7/2}}{35N^{5/2}}+\cdots.
\]
The coefficient $2/15$ combines block weights, pair depletion and the permutation denominator. It is not the raw weight of a single five-element block.
\end{proof}

\subsection{Keeping the denominator exact}

For $2\le k\le N$, define
\begin{equation}\label{eq:globalformula}
\mathcal A_{N,k}
=D_kH_k e^{Nq(\sqrt{K/N})}
(1+\sqrt{K/N})^{\epsilon-1/2}.
\end{equation}
Every factor is positive and finite, including at $k=N$.

\begin{theorem}[Uniform approximation at every order]\label{thm:global}
As $N\to\infty$,
\begin{equation}\label{eq:globalerror}
\sup_{2\le k\le N}\left|\frac{R_k}{\mathcal A_{N,k}}-1\right|
=O(N^{-1/4}).
\end{equation}
Uniformly for $\sqrt N\le k\le N$, the error is
$O(K^{-1}+e^{-\lambda/8})$, and hence $O(N^{-1/2})$.
\end{theorem}
\begin{proof}
Above $\sqrt N$, divide \eqref{eq:normalizedsaddle} by the exact $H_k$. The residual factor is $(1+(-1)^\epsilon e^{-2\lambda})^{-1}$. It differs from one exponentially little. The factor $D_k$ cancels exactly, so no approximation to $(N)_k$ is needed even at $k=N$.

Below $\sqrt N$, compare \eqref{eq:globalformula} to $e^{k^2/(4N)}H_k$. For $0\le x\le1/2$, integrating the elementary bounds on $\log(1+x)$ gives
\[
0\le q(x)+x^4/4\le2x^5/15.
\]
Also $0\le\log D_k-k(k-1)/(2N)\le k^3/(3N^2)$ for $k\le N/4$. The difference between $k(k-1)/(2N)$ and $(k^2+K^2)/(4N)$ is $-k/(2N)$ for even $k$ and $-1/(4N)$ for odd $k$. Thus, for $N\ge16$ and $2\le k\le\sqrt N$,
\[
\left|\log\frac{\mathcal A_{N,k}}{e^{k^2/(4N)}H_k}\right|
\le\frac{19}{30}N^{-1/4}+\frac56N^{-1/2}.
\]
Combine this with the explicit $O(N^{-1/4})$ estimate in Theorem~\ref{thm:rate}. This proves \eqref{eq:globalerror}. The global rate is not asserted to be optimal.
\end{proof}

At the top order, Lemma~\ref{lem:saddle} can be used with degree $K=N$ and offset zero, without a parity restriction on that lemma. Here $r=1/2$ and $b=4N$, giving
\[
[z^N]e^{Nf(z)}=\frac{e^N}{\sqrt{8\pi N}}(1+O(N^{-1})),\qquad
a_N(N)=\frac{1+O(N^{-1})}{2N!}.
\]
This endpoint also follows from the classical truncated-exponential identity for $L_N^{(-N-1)}$. It is a consistency check on the full-range formula, rather than a separate priority claim.
